% ==============================================================================
% 03_yeu_cau.tex — Phần 3. Phân tích yêu cầu (TO-BE)
% Phụ trách: Thành viên 3 (bám sát glossary_lock.md)
% Ngân sách trang: 5 trang
% ==============================================================================

\pttkPart{Phân tích yêu cầu (TO-BE)}

\todosection{Thành viên 3}{Dùng đúng 100\% mã FR-01..FR-22, NFR-01..NFR-10, BR-01..BR-10 trong \texttt{docs/pttk/glossary\_lock.md}. Ngân sách mục tiêu: 5 trang.}

\subsection{Tác nhân (Actors)}
% Bảng tác nhân con người & hệ thống ngoài kèm quyền điển hình.

\subsection{Yêu cầu chức năng}
% Bảng danh sách FR-01..FR-22 theo ưu tiên MoSCoW (dùng longtable).

\subsection{User Story và tiêu chí chấp nhận}
% Viết 3 User Story trọng tâm kèm kịch bản Gherkin Given-When-Then (FR-01, FR-14, FR-17).

\subsection{Yêu cầu phi chức năng}
% Bảng NFR-01..NFR-10 định lượng theo chuẩn ISO/IEC 25010.

\subsection{Quy tắc nghiệp vụ}
% Bảng BR-01..BR-10.

\subsection{Quy trình nghiệp vụ mới (TO-BE)}
\todofig{Hình 2. Quy trình phối hợp lập trình và tự động hóa trên DevFlow (TO-BE) — diagrams/out/hinh\_02\_tobe.pdf}

\subsection{Thuật ngữ (Glossary)}
% Bảng định nghĩa thuật ngữ chuyên ngành DevFlow.
